\section*{Compresión gauge y testigo físico de la brecha}

Sea
\[
 \widehat{\mathscr H}_m
 :=L^2(\widehat X_m,\widehat\mu_m^{\rm ov})
\]
el espacio de la realización material completa. Además de la acción gauge
usual sobre enlaces y marcos, cada collar \(c\) porta la acción finita de
incidencia
\begin{equation}\label{eq:amp-ym-accion-gauge-incidencia}
 q_c\longmapsto q_c+Bx_c,
 \qquad x_c\in\mathbb F_3^4,
\end{equation}
donde las seis filas de \(B\) están indexadas por
\((01,02,03,12,13,23)\). El producto de ambas acciones preserva
\(\widehat\mu_m^{\rm ov}\). Con \(\mathcal C_m\) el conjunto finito de
collares presentes en el nivel, se escribe
\[
 \mathcal G_m^{\rm full}
 :=G^{V(\Lambda_m)}
   \times\prod_{c\in\mathcal C_m}\mathbb F_3^4.
\]
Su promedio de Haar define la proyección
ortogonal
\begin{equation}\label{eq:amp-ym-proyeccion-fisica-haar}
 \mathbb P_m^{\rm phys}F
 :=\int_{\mathcal G_m^{\rm full}}F(g^{-1}\!\cdot\,)\,dg,
 \qquad
 \widehat{\mathscr H}_m^{\rm phys}
 :=\operatorname{Ran}\mathbb P_m^{\rm phys}.
\end{equation}
Ésta es una compresión sobre la subálgebra gauge del estado completo; la
marginal que olvida las coordenadas de incidencia es otro mapa y no se
confunde con \(\mathbb P_m^{\rm phys}\).

\begin{lemma}[Reducción del Hamiltoniano a la subálgebra física]
\label{lem:amp-ym-compresion-fisica}
La proyección \(\mathbb P_m^{\rm phys}\) reduce el generador y su
semigrupo. Las inclusiones de refinamiento reducen simultáneamente las
subálgebras físicas:
\begin{equation}\label{eq:amp-ym-compresion-entrelazada}
 \begin{gathered}
 [\mathbb P_m^{\rm phys},\widehat H_m^{\rm ov}]=0,
 \qquad
 [\mathbb P_m^{\rm phys},e^{-t\widehat H_m^{\rm ov}}]=0,\\
 \mathbb P_{m+1}^{\rm phys}\widehat J_m
 =\widehat J_m\mathbb P_m^{\rm phys},\qquad
 H_m^{\rm phys}
 :=\widehat H_m^{\rm ov}\big|_{\widehat{\mathscr H}_m^{\rm phys}},\\
 H_{m+1}^{\rm phys}\widehat J_m
 =\widehat J_mH_m^{\rm phys}.
 \end{gathered}
\end{equation}
\end{lemma}

\begin{proof}
El generador físico sobre enlaces es gauge equivariante por la
biinvariancia de Haar. En la carta producto, \(E_{q_c}\) es la esperanza
uniforme en \(\mathbb F_3^6\); por invariancia de la medida uniforme bajo
las traslaciones de \eqref{eq:amp-ym-accion-gauge-incidencia}, conmuta con
esa acción. Las expectativas de las subcolas de marcación actúan en
factores disjuntos. Por tanto cada sumando de
\[
 \widehat H_m^{\rm ov}
 =H_m^{\rm ov}\otimes I
 +3\kappa_{\rm av}\sum_c
 \bigl[(I-E_{q_c})+(I-E_{u_c^{\rm P}})\bigr]
\]
conmuta con el promedio \eqref{eq:amp-ym-proyeccion-fisica-haar}. El
cálculo funcional da la conmutación con el semigrupo. Finalmente,
\(\widehat J_m\) conserva los factores ancestrales e inserta la constante
en los factores nuevos; promediar antes o después de esa inclusión produce
la misma función. Esto prueba las cuatro identidades de
\eqref{eq:amp-ym-compresion-entrelazada}.
\end{proof}

\begin{lemma}[Testigo gauge no nulo y entrelazado]
\label{lem:amp-ym-testigo-fisico-q6}
Para todo collar \(c\), la función
\begin{equation}\label{eq:amp-ym-testigo-fisico-def}
 W_c(q)
 :=\mathbf 1_{\{q_1+\cdots+q_6=0\}}-\frac13
\end{equation}
pertenece a \(\widehat{\mathscr H}_m^{\rm phys}\), no es nula y verifica
\begin{equation}\label{eq:amp-ym-testigo-fisico-momentos}
 \mathbb E W_c=0,\qquad
 \|W_c\|^2=\frac29,\qquad
 \mathbb E(W_c^3)=\frac2{27}.
\end{equation}
Además,
\begin{equation}\label{eq:amp-ym-testigo-fisico-gap}
 H_m^{\rm phys}W_c=3\kappa_{\rm av}W_c,\qquad
 H_{m+1}^{\rm phys}\widehat J_mW_c
 =3\kappa_{\rm av}\widehat J_mW_c.
\end{equation}
En consecuencia, el valor \(3\kappa_{\rm av}\) pertenece al espectro de la
restricción física y no sólo al de una dilatación material.
\end{lemma}

\begin{proof}
La acción gauge sobre enlaces actúa en el primer factor y deja fijo
\(W_c\). Para la acción de incidencia, cada \(x_i\) aparece en tres de las
seis coordenadas, luego
\[
 \sum_{i<j}(Bx)_{ij}=3\sum_i x_i=0
 \quad\text{en }\mathbb F_3.
\]
La condición de \eqref{eq:amp-ym-testigo-fisico-def} es, por tanto,
invariante; también lo es bajo las permutaciones de las seis incidencias.
Así \(\mathbb P_m^{\rm phys}W_c=W_c\).

Bajo la medida uniforme, la suma de los seis trits es uniforme en
\(\mathbb F_3\). La función toma el valor \(2/3\) con probabilidad \(1/3\)
y \(-1/3\) con probabilidad \(2/3\), lo que da
\eqref{eq:amp-ym-testigo-fisico-momentos}. El primer sumando de
\(\widehat H_m^{\rm ov}\) anula \(W_c\); para \(c'\ne c\),
\((I-E_{q_{c'}})W_c=0\), mientras
\(E_{q_c}W_c=0\). Todas las expectativas de marcación anulan la
diferencia correspondiente porque \(W_c\) es constante en esas subcolas.
Queda exactamente
\(\widehat H_m^{\rm ov}W_c=3\kappa_{\rm av}W_c\). La primera identidad de
\eqref{eq:amp-ym-testigo-fisico-gap} sigue de la reducción; la segunda,
del entrelazamiento de \eqref{eq:amp-ym-compresion-entrelazada}.
\end{proof}

\section*{Semigrupo propietario, vacío y cota espectral exacta}

Sobre la subálgebra física se escribe
\[
 D_m^{\rm YM}:=H_m^{\rm phys},\qquad
 K_{m,t}:=\exp(-tD_m^{\rm YM}),\qquad
 P_{\Omega_m}:=|\Omega_m\rangle\langle\Omega_m|.
\]
La descomposición de Hoeffding del mismo generador y el cálculo funcional
dan, sin cambiar de medida ni de proceso,
\begin{equation}\label{eq:amp-ym-owner-semigrupo-vacio}
 K_{m,t}P_{\Omega_m}=P_{\Omega_m},
 \qquad
 \bigl\|K_{m,t}(I-P_{\Omega_m})\bigr\|
 \leq e^{-3\kappa_{\rm av}t}.
\end{equation}
El testigo de \eqref{eq:amp-ym-testigo-fisico-def} es no nulo y alcanza la
cota:
\begin{equation}\label{eq:amp-ym-owner-testigo-wc}
 D_m^{\rm YM}W_c=3\kappa_{\rm av}W_c,
 \qquad \|W_c\|^2=\frac29.
\end{equation}
Por tanto el vacío es simple y la brecha no sólo está acotada inferiormente,
sino que es exactamente
\begin{equation}\label{eq:amp-ym-owner-brecha-exacta}
 \Delta_{\rm YM}^{\rm HMT}=3\kappa_{\rm av}>0.
\end{equation}
El semigrupo, el proyector de vacío, la contracción del complemento y
\(W_c\) pertenecen a la misma familia proyectiva y constituyen la cadena
propietaria del cierre.

\section*{Reconstrucción euclídea, campos de Schwinger y lector de curvatura}

La medida proyectiva común del teorema determina inclusiones isométricas
\(\widehat J_m\) entre los espacios de nivel. Para cada una de las cuatro
reflexiones coordenadas \(\theta_\mu\), \(\mu=1,\ldots,4\), existe una
factorización
\begin{equation}\label{eq:amp-ym-cuatro-os}
 \int\overline{F(\theta_\mu y)}F(y)\,
 d\widehat\Omega_m(y)
 =\|\mathcal B_{m,\mu}F\|_2^2\geq0.
\end{equation}
Las cuatro formas proceden de la misma medida y del mismo semigrupo. Los
conectores OS
\[
 \mathcal J_{m,\mu}^{\rm OS}
 =\mathcal V_{m+1,\mu}^{-1}\widehat J_m\mathcal V_{m,\mu}
\]
entrelazan sus Hamiltonianos. En el colímite,
\begin{equation}\label{eq:amp-ym-hamiltoniano-comun}
 \mathcal V_{\infty,\mu}H_{{\rm OS},\infty}^{(\mu)}
 \mathcal V_{\infty,\mu}^{-1}
 =\widehat H_\infty,
 \qquad \mu=1,\ldots,4.
\end{equation}

La forma límite se diagonaliza por la descomposición de Hoeffding:
\begin{equation}\label{eq:amp-ym-forma-limite}
 \mathcal E_\infty(u)=3\kappa_{\rm av}
 \sum_{S\Subset\mathcal I_\infty}|S|\,\|u_S\|^2.
\end{equation}
Las sumas parciales proporcionan una sucesión de recuperación y la
semicontinuidad produce el límite inferior. La convergencia de Mosco de las
formas cerradas conserva el núcleo unidimensional y el umbral espectral. El
vector material \(W_c\) satisface
\begin{equation}\label{eq:amp-ym-testigo-gap}
 \widehat H_mW_c=3\kappa_{\rm av}W_c,
 \qquad \|W_c\|^2=\frac29,
 \qquad \mathbb E(W_c^3)=\frac2{27},
\end{equation}
de modo que la cota inferior se alcanza y la brecha es exactamente
\(3\kappa_{\rm av}\).

\begin{theorem}[Terminación euclídea y publicación de curvatura]
\label{thm:amp-ym-terminacion}
La familia proyectiva del teorema \ref{thm:realizaciones-ym} determina una
acción fuertemente continua de \(E(4)\), una red local gauge no escalar,
funcionales de Schwinger compatibles y una publicación de campo en
\(H^{-s}_{\rm loc}(\mathbb R^4)\) para \(s>2\). El producto superficial
ordenado de los enlaces define lectores clover y \(F^2\); en el dominio
suave, el lector de acción converge a
\begin{equation}\label{eq:amp-ym-f2}
 \mathcal S_{\rm YM}(A)
 =\frac1{4g_R^2}\int_{\mathbb R^4}\|F_A(x)\|^2\,d^4x.
\end{equation}
Estas publicaciones conservan el vacío simple y la brecha exacta del mismo
Hamiltoniano común.
\end{theorem}

\begin{proof}
Los Gram cofinales de los lectores suavizados y la conmutatividad de los
cuadrados de rotación y refinamiento preservan los ideales nulos; su
colímite construye la representación fuerte de \(E(4)\) y la red local.
Las estimaciones martingala de los incrementos de enlace dan tightness en
\(H^{-s}_{\rm loc}\), \(s>2\), y las características mixtas identifican los
funcionales de Schwinger con el mismo proceso cilíndrico. El producto
orientado alrededor de cada plaqueta construye el clover. Su expansión en
el régimen suave y la convergencia fuerte del lector marginal dan
\eqref{eq:amp-ym-f2}. Ninguna de estas operaciones modifica la medida, el
semigrupo ni la forma límite usados para obtener
\eqref{eq:amp-ym-testigo-gap}; por ello el vacío y la brecha se conservan.
\end{proof}

\section*{Identificación exacta del límite gauge}

La identificación de la teoría no descansa en que cuatro reflexiones
generen por sí solas el grupo euclídeo.  Las reflexiones proporcionan la
positividad OS; la acción continua procede, separadamente, del grupoide de
pantallas y de la compatibilidad de refinamiento.  Para
\(g=(g_v)_{v\in\Lambda_m}\in G^{\Lambda_m}\), la acción sobre una variable
de enlace es
\begin{equation}\label{eq:amp-ym-accion-gauge-enlace}
 (g\cdot U)_e=g_{s(e)}U_e g_{t(e)}^{-1}.
\end{equation}
La invariancia de Haar, la desintegración triangular del kernel y la
cancelación de las aristas interiores prueban
\begin{equation}\label{eq:amp-ym-invariancia-gauge-medida}
 (g\cdot)_*\widehat\mu_m^{\rm ov}=\widehat\mu_m^{\rm ov},
 \qquad
 K_{m,t}^{\rm ov}(g\cdot U,g\cdot dV)
 =K_{m,t}^{\rm ov}(U,dV).
\end{equation}
Por diferenciación sobre el núcleo cilíndrico, todo generador vertical
\(X\) satisface la identidad de Ward
\begin{equation}\label{eq:amp-ym-ward}
 \int \nabla_XF\,d\widehat\mu_m^{\rm ov}=0.
\end{equation}

La curvatura se obtiene del mismo sistema de enlaces.  Si una plaqueta
gruesa \(p\) se subdivide en sus 81 subplaquetas \(p'\), entonces
\begin{equation}\label{eq:amp-ym-producto-superficial}
 \operatorname{Hol}^{(m)}_{p}
 =\prod_{p'\subset p}^{\longrightarrow}
 T_{p'}\operatorname{Hol}^{(m+1)}_{p'}T_{p'}^{-1}.
\end{equation}
Las contribuciones de toda arista interior aparecen con orientaciones
opuestas y se cancelan exactamente.  La expresión es gauge-covariante y
conmuta con las inclusiones \(\widehat J_m\).  En una carta suave se define
\begin{equation}\label{eq:amp-ym-clover-definido}
 F_{{\rm cl},m}(x)
 =\frac1{4a_m^2}\sum_{p\ni x}
 \operatorname{Ad}_{T_p}\log\operatorname{Hol}_{p},
 \qquad
 F_{{\rm cl},m}=F_A+O(a_m^2).
\end{equation}
De aquí, para \(A\in C^4_{\rm loc}\) y
\(a_m^2/g_m^2\to0\), se obtiene por suma de Riemann
\begin{equation}\label{eq:amp-ym-f2-convergencia-calculada}
 \frac{a_m^4}{4g_m^2}
 \sum_{x\in\Lambda_m}\|F_{{\rm cl},m}(x)\|^2
 \longrightarrow
 \frac1{4g_R^2}\int_{\mathbb R^4}\|F_A(x)\|^2\,d^4x.
\end{equation}

\begin{theorem}[Proceso gauge euclídeo y lector de curvatura de
Yang--Mills]
\label{thm:amp-ym-identificacion-gauge}
El límite construido en el teorema
\ref{thm:realizaciones-ym} es un proceso gauge euclídeo en dimensión
\(4\) para el grupo compacto, conectado y simple fijado. Posee
covariancia gauge local, una red local no escalar, positividad por
reflexión, una representación fuertemente continua de \(E(4)\),
funcionales de Schwinger compatibles, curvatura gauge-covariante con límite
clásico \eqref{eq:amp-ym-f2-convergencia-calculada}, vacío simple y brecha
espectral
\(3\kappa_{\rm av}>0\).
\end{theorem}

\begin{proof}
Las identidades \eqref{eq:amp-ym-invariancia-gauge-medida}--
\eqref{eq:amp-ym-ward} construyen la simetría gauge y sus identidades de
Ward en todos los niveles; su naturalidad las transporta al límite.  La
localidad uniforme de los kernels produce la red isótona y local.  Las
cuatro formas \eqref{eq:amp-ym-cuatro-os}, junto con la acción del grupoide
de pantallas sobre el núcleo Weyl, producen respectivamente la positividad
OS y la acción fuerte de \(E(4)\).  La tightness en
\(H^{-s}_{\rm loc}\), \(s>2\), identifica una única ley límite y sus
funcionales de Schwinger.  Finalmente,
\eqref{eq:amp-ym-producto-superficial}--
\eqref{eq:amp-ym-f2-convergencia-calculada} identifican el lector de campo
y su límite clásico, mientras
\eqref{eq:amp-ym-forma-limite}--
\eqref{eq:amp-ym-testigo-gap} prueban vacío simple y brecha exacta para la
misma ley.

La medida queda determinada constructivamente por sus marginales
proyectivas, su kernel reversible y sus funcionales cilíndricos. El lector
clover no repondera esa medida: publica una función del mismo proceso y su
límite clásico.
\end{proof}

\begin{proposition}[Control no gobernante de comparación Gibbs--acción]
\label{prop:amp-ym-puente-accion-requerido}
La ley proyectiva, su vacío simple y su brecha exacta ya están construidos
por \eqref{eq:amp-ym-owner-semigrupo-vacio}--
\eqref{eq:amp-ym-owner-brecha-exacta}. Como comparación exterior opcional de
tipo Gibbs--acción, una parametrización por \(g_R\) puede estudiar una familia
\[
 g_R\longmapsto
 \bigl(\widehat\mu_{m,g_R}^{\rm ov},H_{m,g_R}^{\rm phys}\bigr)
\]
y una relación entre la medida o su forma de Dirichlet y el funcional de
acción, por ejemplo
\begin{equation}\label{eq:amp-ym-identidad-accion-requerida}
 \frac{d\widehat\mu_{m,g_R}^{\rm ov}}{d\nu_m}(U)
 =Z_{m,g_R}^{-1}e^{-S_{m,g_R}^{F^2}(U)},
 \qquad
 \mathfrak h_{m,g_R}[F]
 =\langle F,H_{m,g_R}^{\rm phys}F\rangle,
\end{equation}
con compatibilidad de refinamiento y paso al límite. Aquí \(\nu_m\) puede
ser la medida producto de Haar de la red finita. Esta comprobación se rotula
\texttt{CONTROL\_NO\_GOBERNANTE}: no es premisa de la teoría gauge HMT ni
de su brecha.
\end{proposition}

\begin{proof}
En \eqref{eq:amp-ym-f2-convergencia-calculada}, variar \(g_R\) reescala el
lector de curvatura, mientras las fórmulas que construyen
\(\widehat\mu_m^{\rm ov}\), \(\widehat H_m^{\rm ov}\) y la brecha
\(3\kappa_{\rm av}\) no contienen \(g_R\). Son, por ello, construcciones
independientes. La compresión
\eqref{eq:amp-ym-compresion-entrelazada} y el testigo
\eqref{eq:amp-ym-testigo-fisico-gap} prueban la persistencia física de la
brecha del proceso ya construido; el control sólo compara después una
dependencia dinámica adicional en \(g_R\).
\end{proof}

La secuencia demostrada en esta sección es
\[
 \text{medida común}\to\text{cuatro OS}\to\text{Hamiltoniano común}
 \to E(4),\ \text{Schwinger},\ H^{-s}_{\rm loc}
 \to\text{clover y }F^2.
\]
Esta secuencia demuestra una teoría Yang--Mills cuatridimensional con vacío
simple y brecha positiva exacta \(3\kappa_{\rm av}\) en el dominio HMT
registrado. La identidad
\eqref{eq:amp-ym-identidad-accion-requerida} queda fuera de la cadena
rectora como \texttt{CONTROL\_NO\_GOBERNANTE}.
